\chapter{Wieże Postnikowa, iloczyny Whiteheada i operacje Steenroda}
\label{ch:postnikow-v7}

\section{Dlaczego grupy homotopii nie wystarczają?}
Niech $X$ będzie łukowo spójnym CW-kompleksem. Grupy
$\pi_n(X)$ opisują odwzorowania sfer, ale zapisanie samej
rodziny tych grup nie mówi jeszcze, jak odwzorowania
różnych wymiarów oddziałują ze sobą. Właśnie tę brakującą
informację zachowują \emph{niezmienniki Postnikowa}. Z kolei
\emph{iloczyny Whiteheada} mierzą oddziaływanie elementów
grup homotopii, a \emph{operacje Steenroda} wykrywają
informację niewidoczną w addytywnych grupach kohomologii
czy nawet w samych kwadratach iloczynu kubkowego.

Podamy najpierw konstrukcje dla \emph{jednospójnych}
CW-kompleksów z wybranym punktem bazowym. Jednospójność
oznacza tu $\pi_0(X)=*$ i $\pi_1(X)=0$. Dzięki temu
na wyższych grupach homotopii nie występuje nietrywialne
działanie grupy podstawowej. Dopiero w ostatniej sekcji
wyjaśnimy, co zmienia się w przypadku ogólnym.

\begin{definicja}
Niech $n\ge1$. \emph{$n$-typem homotopijnym} nazywamy
łukowo spójny CW-kompleks $Y$, dla którego $\pi_j(Y)=0$
dla każdego $j>n$. Odwzorowanie $f:X\to Y$ jest
\emph{$n$-równoważnością} w tym rozdziale, jeśli
$f_*:\pi_j(X)\to\pi_j(Y)$ jest izomorfizmem dla
$1\le j\le n$ (przy zgodnym wyborze punktów bazowych).
\emph{$n$-etapem Postnikowa} przestrzeni $X$ nazywamy
parę $(P_nX,p_n)$, gdzie $P_nX$ jest $n$-typem,
a $p_n:X\to P_nX$ jest $n$-równoważnością.
\end{definicja}

\begin{uwaga}
Zdefiniowane tu pojęcie $n$-równoważności jest słabsze
od spotykanego w literaturze określenia \emph{$n$-spójne odwzorowanie},
które dodatkowo narzuca warunek surjektywności w wymiarze
$n+1$ lub warunek na włókno homotopijne, zależnie od
konwencji. Nie będziemy mieszać tych terminów.
\end{uwaga}

\begin{twierdzenie}[Istnienie etapów Postnikowa]
\label{thm:postnikow-istence-v7}
Dla każdego łukowo spójnego CW-kompleksu $X$ i każdej
liczby $n\ge1$ istnieje etap $(P_nX,p_n)$.
Można dobrać etapy w wieżę odwzorowań
$\cdots\to P_{n+1}X\to P_nX\to\cdots\to P_1X$
zgodnych z odwzorowaniami $p_n$ w kategorii homotopijnej.
\end{twierdzenie}
\begin{proof}
Do $X$ dołączamy komórki wymiaru co najmniej $n+2$.
Na początku wybieramy reprezentantów wszystkich elementów
$\pi_{n+1}(X)$ i dołączamy do ich odwzorowań
$S^{n+1}\to X$ komórki wymiaru $n+2$.
Długi ciąg homotopii pary oraz względne twierdzenie
o przybliżeniu komórkowym pokazują, że nowe komórki
nie zmieniają $\pi_j$ dla $j\le n$, natomiast zabijają
wybrane elementy $\pi_{n+1}$. Jeżeli na tym etapie
pozostałyby elementy $\pi_{n+1}$, dołączamy komórki
wzdłuż ich reprezentantów; równoważnie można od razu
zabrać reprezentantów całej grupy. Następnie zabijamy
$\pi_{n+2}$ komórkami wymiaru $n+3$, i tak dalej.

Niech $Y$ będzie sumą kolejnych podkompleksów w topologii
CW. Dla ustalonego $j>n$ wszystkie elementy $\pi_j$
któregoś etapu giną po dołączeniu komórek $(j+1)$-wymiarowych,
a późniejsze komórki o wymiarze co najmniej $j+2$
nie odtwarzają $\pi_j$. Każda odwzorowanie $S^j\to Y$ jest
homotopijna do odwzorowania trafiającej do skończonego podkompleksu
$Y$, a więc do pewnego etapu konstrukcji (twierdzenie
o zwartościowym charakterze odwzorowań do CW-kompleksu).
Otrzymujemy $\pi_j(Y)=0$. Ponieważ dołączaliśmy tylko
komórki wymiaru co najmniej $n+2$, dla $j\le n$
odwzorowanie $X\to Y$ indukuje izomorfizm na $\pi_j$.
Przyjmujemy $P_nX=Y$.

Aby uzyskać odwzorowania $P_{n+1}X\to P_nX$, przedłużamy
$p_n:X\to P_nX$ do $P_{n+1}X$ komórka po komórce.
Komórki dołączone do $X$ w konstrukcji $P_{n+1}X$
mają wymiary co najmniej $n+3$, a zatem przeszkody
przedłużenia leżą w grupach $\pi_j(P_nX)$ z $j\ge n+2$,
które są zerowe. Złożenie z $p_{n+1}$ jest homotopijne
z $p_n$. Wybory przedłużeń nie są kanoniczne na poziomie
pojedynczych przestrzeni, ale wyznaczają poprawną wieżę
w kategorii homotopijnej.
\end{proof}

\begin{przyklad}[Dwa pierwsze etapy sfery]
Dla $X=S^2$ mamy $P_1S^2\simeq *$, a
$P_2S^2\simeq K(\Z,2)\simeq\C P^\infty$.
Nie wynika stąd, że $S^2\simeq\C P^\infty$:
pierwsza przestrzeń ma niezerową grupę $\pi_3\cong\Z$,
a druga już nie. Konstrukcja $P_3S^2$ musi więc
zawierać informację, której nie ma w samym $K(\Z,2)$.
\end{przyklad}

\section{Włókna kolejnych etapów i niezmienniki \texorpdfstring{$k$}{k}}
\begin{definicja}[Włókno homotopijne]
Niech $f:B\to Y$ będzie odwzorowaniem przestrzeni z punktami
bazowymi i $y_0\in Y$. \emph{Włóknem homotopijnym} $f$
nad $y_0$ nazywamy przestrzeń
\[
\operatorname{hofib}(f)=\{(b,\gamma):b\in B,
\ \gamma:[0,1]\to Y,\ \gamma(0)=f(b),\ \gamma(1)=y_0\},
\]
z topologią podprzestrzeni $B\times Y^{[0,1]}$,
gdzie $Y^{[0,1]}$ ma topologię zwarto-otwartą.
\end{definicja}

\begin{lemat}[Włókno odwzorowania do przestrzeni Eilenberga--Mac Lane'a]
\label{lem:hofib-K-v7}
Dla $n\ge1$ i grupy abelowej $A$ włókno homotopijne
odwzorowania $k:B\to K(A,n+1)$ jest włóknem fibracji
nad $B$ o typie homotopijnym $K(A,n)$.
Dla spójnej $B$ działanie monodromii na $A$ jest trywialne.
\end{lemat}
\begin{proof}
Rzutowanie $\operatorname{hofib}(k)\to B$ jest cofnięciem
fibracji dróg $P_{y_0}K(A,n+1)\to K(A,n+1)$,
gdzie $P_{y_0}Y$ składa się z dróg kończących się w $y_0$.
Włókno jest przestrzenią pętli $\Omega K(A,n+1)$.
Z $\pi_j(\Omega Y)=\pi_{j+1}(Y)$ wynika, że ma ona
jedyną niezerową grupę homotopii $\pi_n=A$,
więc jest słabo równoważna $K(A,n)$; po wybraniu
modeli CW daje to równoważność homotopijną.
Działanie $\pi_1(B)$ otrzymujemy przez cofnięcie
działania $\pi_1(K(A,n+1))$, a ta grupa jest zerowa
dla $n\ge1$.
\end{proof}

\begin{twierdzenie}[Rozkład Postnikowa dla przestrzeni jednospójnej]
\label{thm:postnikov-k-v7}
Niech $X$ będzie jednospójnym CW-kompleksem, $n\ge2$,
a $A_n=\pi_n(X)$ (grupa abelowa). Etapy Postnikowa
można dobrać wraz z fibracjami homotopijnymi
\[
K(A_n,n)\longrightarrow P_nX\xrightarrow{q_n}P_{n-1}X.
\]
Fibracja ta jest, z dokładnością do równoważności
włóknowej, włóknem homotopijnym pewnej odwzorowania
\(k_n:P_{n-1}X\to K(A_n,n+1)\).
Klasa $[k_n]\in H^{n+1}(P_{n-1}X;A_n)$
nazywa się \emph{$n$-tym niezmiennikiem Postnikowa}.
Zależy od wyboru identyfikacji grup homotopii,
a bez tych identyfikacji jest określona z dokładnością
do działania odpowiednich automorfizmów.
\end{twierdzenie}
\begin{proof}
Rozpatrzmy odwzorowanie $q_n:P_nX\to P_{n-1}X$,
odpowiadającą zgodnym etapom. Zastępujemy ją fibracją
przestrzeni dróg, nie zmieniając typu homotopijnego
przestrzeni całkowitej. Z długiego ciągu homotopii
wynika, że jej włókno homotopijne $F_n$ ma
$\pi_j(F_n)=0$ dla $j\ne n$ oraz
$\pi_n(F_n)\cong\pi_n(X)=A_n$: dla $j<n$
odwzorowanie $q_n$ indukuje izomorfizm, a dla $j>n$
obie grupy w wieży są zerowe. Zatem
$F_n\simeq K(A_n,n)$.

Wyjaśnienia wymaga, dlaczego taka fibracja jest określona
przez jedną klasę kohomologii. Na kolejnych komórkach
jednospójnej podstawy $B=P_{n-1}X$ wybieramy lokalne
podniesienia do $E=P_nX$. Ponieważ włókno
$K(A_n,n)$ ma zerowe grupy $\pi_j$ dla $j<n$,
wybory można uzgodnić do wymiaru $n$. Na
$(n+1)$-komórkach pojawiają się elementy $A_n$;
jak w teorii przeszkód tworzą one kocykl i klasę
$k\in H^{n+1}(B;A_n)$. Wyższe przeszkody
na zgodność wyborów nie występują, ponieważ
$\pi_j(K(A_n,n))=0$ dla $j>n$.

Jest to podstawowy przypadek klasyfikacji \emph{głównych
fibracji $K(A,n)$}. Można go również zobaczyć
bezpośrednio: wybieramy odwzorowanie $k:B\to K(A_n,n+1)$
reprezentującą klasę i konstruujemy jej włókno
homotopijne zgodnie z Lematem~\ref{lem:hofib-K-v7}.
Na każdym szkielecie obie fibracje mają ten sam
kołańcuch przeszkody, a indukcja komórkowa daje
ich równoważność włóknową. Ważna jest tu \emph{prostota}
struktury homotopijnej: $\pi_1(B)=0$ i dlatego
nie występują skręcone współczynniki ani niezależne
działania monodromii na $A_n$.
\end{proof}

\begin{uwaga}[Granica powyższego dowodu]
Identyfikacja włókna przez długi ciąg homotopii jest
pełnym rachunkiem. Natomiast klasyfikacja wszystkich
głównych fibracji $K(A,n)$ przez $H^{n+1}$ korzysta
z twierdzenia o klasyfikacji takich fibracji, którego
konstrukcję komórkową podaliśmy w argumentacji.
Nie twierdzimy, że \emph{dowolna} fibracja o włóknie
$K(A,n)$ jest główna: bez dodatkowych danych może
wystąpić nietrywialne działanie grupy automorfizmów
włókna. Wieża Postnikowa przestrzeni jednospójnej
ma wymagany typ główny.
\end{uwaga}

\begin{center}
\begin{tikzpicture}[>=Stealth, every node/.style={font=\small}]
\node (X) at (-3.6,1.8) {$X$};
\node (P4) at (-1.5,1.8) {$P_4X$};
\node (P3) at (0.7,1.8) {$P_3X$};
\node (P2) at (3.0,1.8) {$P_2X$};
\draw[->,thick] (X)--(P4) node[midway,above]{$p_4$};
\draw[->,thick] (P4)--(P3) node[midway,above]{$q_4$};
\draw[->,thick] (P3)--(P2) node[midway,above]{$q_3$};
\node (K4) at (-1.5,.35) {$K(\pi_4X,4)$};
\node (K3) at (.7,.35) {$K(\pi_3X,3)$};
\draw[->,thick] (K4)--(P4);
\draw[->,thick] (K3)--(P3);
\node (K5) at (-1.5,-1.0) {$K(\pi_4X,5)$};
\node (K4b) at (.7,-1.0) {$K(\pi_3X,4)$};
\draw[->,thick,dashed] (P3) to[out=-125,in=10] (K5)
 node[midway,left] {$k_4$};
\draw[->,thick,dashed] (P2) to[out=-130,in=10] (K4b)
 node[midway,right] {$k_3$};
\end{tikzpicture}
\end{center}
\noindent Schemat wieży: włókno odwzorowania $q_n$ ma typ $K(\pi_nX,n)$,
a klasa $k_n$ jest reprezentowana przez odwzorowanie z $P_{n-1}X$
do $K(\pi_nX,n+1)$. Przerywane strzałki oznaczają
odwzorowania klasyfikujące, a nie części ciągu fibracji.

\section{Niezmiennik \texorpdfstring{$k_3$}{k3} dla \texorpdfstring{$S^2$}{S2}}
\begin{twierdzenie}[Pierwszy nietrywialny etap sfery]
\label{thm:s2-k-v7}
Niech $u\in H^2(K(\Z,2);\Z)$ będzie generatorem,
wybranym tak, że $H^*(K(\Z,2);\Z)\cong\Z[u]$.
Dla $X=S^2$ niezmiennik $k_3$ odpowiada klasie
$\pm u^2\in H^4(K(\Z,2);\Z)\cong\Z$.
W szczególności $P_3S^2$ jest włóknem homotopijnym
odwzorowania reprezentującej $\pm u^2$.
\end{twierdzenie}
\begin{proof}
Mamy $\pi_2(S^2)=\Z$ i $\pi_3(S^2)=\Z$,
stąd $P_2S^2\simeq K(\Z,2)$, zaś włókno
$P_3S^2\to P_2S^2$ ma typ $K(\Z,3)$.
Kohomologiczny ciąg spektralny Serre'a tej fibracji
z całkowitymi współczynnikami ma w niskich stopniach
$E_2^{p,q}=H^p(K(\Z,2);H^q(K(\Z,3);\Z))$.
Dla $p+q\le4$ wiersze $q=0$ i $q=3$
są jedynymi mogącymi uczestniczyć w różniczce
zaczynającej się w $E_4^{0,3}=\Z$;
jej celem jest $E_4^{4,0}=\Z\cdot u^2$.

Wybierzmy model $P_3S^2$ skonstruowany z $S^2$
przez dołączanie komórek wymiaru co najmniej $5$.
Nie ma wtedy nowych względnych komórek w stopniu $4$,
więc
$H^4(P_3S^2,S^2;\Z)=0$ i
$H^5(P_3S^2,S^2;\Z)$ może być niezerowa.
Z długiego ciągu pary wynika, że
$H^4(P_3S^2;\Z)=0$, ponieważ
$H^4(S^2;\Z)=0$ i $H^4(P_3S^2,S^2;\Z)=0$.
W konsekwencji klasa $u^2$ musi zginąć przed
stroną $E_\infty$, a jedyną różniczką mogącą
ją zabić jest $d_4:E_4^{0,3}\to E_4^{4,0}$.
Jej obraz musi być całym $\Z$, zatem jest to
mnożenie przez $\pm1$. Z definicji transgresji
klasa $k_3$ jest obrazem generatora włókna
przez tę różniczkę, co daje $k_3=\pm u^2$.
\end{proof}

\begin{uwaga}
W dowodzie wystarczy silniejszy fakt komórkowy:
$P_3S^2$ uzyskujemy z $S^2$ przez dołączanie komórek
w wymiarach co najmniej $5$. Nie wynika z tego
izomorfizm $H^5$, ani ogólne izomorfizmy wyższej
kohomologii. Informacja o \emph{pierwszej} niezerowej
różniczce w ciągu Serre'a jest tym, co rozróżnia
$S^2$ od $K(\Z,2)$.
\end{uwaga}


\begin{przyklad}[\texorpdfstring{$S^2$}{S2} i
\texorpdfstring{$\C P^2$}{CP2}: jednakowe $\pi_2$,
ale różne $k_3$]
Obie przestrzenie są jednospójne i mają
$\pi_2\cong\Z$, więc ich drugi etap Postnikowa
to $K(\Z,2)$. Natomiast z fibracji Hopfa
$S^1\to S^5\to\C P^2$ wynika
$\pi_3(\C P^2)\cong\pi_3(S^5)=0$.
Zatem $P_3\C P^2\simeq K(\Z,2)$, a odpowiadający
niezmiennik $k_3$ znika, inaczej niż dla $S^2$.
Same drugie grupy homotopii nie odróżniają więc tych
przestrzeni; pierwszy niezerowy niezmiennik Postnikowa
dostarcza dodatkowej informacji.
\end{przyklad}

\section{Iloczyny Whiteheada i komórka górna iloczynu sfer}
\begin{definicja}[Iloczyn Whiteheada]
Niech $X$ będzie jednospójną przestrzenią z punktem
bazowym, $p,q\ge2$, a
$\alpha\in\pi_p(X)$, $\beta\in\pi_q(X)$.
Oznaczmy przez $\iota_p:S^p\to S^p\vee S^q$ i
$\iota_q:S^q\to S^p\vee S^q$ włączenia składników.
Produkt $S^p\times S^q$ ma strukturę CW, w której
$(S^p\times\{*\})\cup(\{*\}\times S^q)=S^p\vee S^q$
jest szkieletem bez ostatniej komórki wymiaru $p+q$.
Niech $\omega:S^{p+q-1}\to S^p\vee S^q$
będzie odwzorowaniem dołączającym tej komórki, z określoną
konwencją orientacji. \emph{Iloczynem Whiteheada}
nazywamy
$[\alpha,\beta]=(\alpha\vee\beta)_*[\omega]
\in\pi_{p+q-1}(X)$, gdzie
$\alpha\vee\beta:S^p\vee S^q\to X$.
\end{definicja}

\begin{uwaga}
Zmiana konwencji orientacji może zmienić globalny znak
iloczynu. W dalszych rachunkach, w których znak ma
znaczenie, podajemy wynik z dokładnością do znaku.
Nie mylimy iloczynu Whiteheada z komutatorem
zwykłej grupy podstawowej: dla $p,q\ge2$ wynik
należy do grupy w stopniu $p+q-1$.
\end{uwaga}

\begin{twierdzenie}[Przeszkoda do przedłużenia na iloczyn]
\label{thm:whitehead-product-v7}
Niech $p,q\ge2$, $X$ będzie przestrzenią z punktem
bazowym, a $\alpha:S^p\to X$ i $\beta:S^q\to X$
odwzorowaniami punktowanymi. Odwzorowanie
$\alpha\vee\beta:S^p\vee S^q\to X$ przedłuża się
na $S^p\times S^q$ wtedy i tylko wtedy, gdy
$[\alpha,\beta]=0$ w $\pi_{p+q-1}(X)$.
\end{twierdzenie}
\begin{proof}
Przestrzeń $S^p\times S^q$ powstaje z bukietu
$S^p\vee S^q$ przez dołączenie jednej komórki
$D^{p+q}$ wzdłuż odwzorowania $\omega$.
Przedłużenie odwzorowania $f=\alpha\vee\beta$
przez tę komórkę jest równoważne przedłużeniu
$f\circ\omega:S^{p+q-1}\to X$ na dysk.
Taka odwzorowanie przedłuża się na dysk dokładnie wtedy,
gdy jest homotopijna do odwzorowania stałego (z uwzględnieniem
punktu bazowego), czyli gdy reprezentowana klasa
$[\alpha,\beta]$ jest zerowa.
\end{proof}

\begin{definicja}[Niezmiennik Hopfa przez komórkę dołączaną]
Niech $n\ge2$ będzie parzyste i
$f:S^{2n-1}\to S^n$ będzie odwzorowaniem punktowanym.
Niech $C_f=S^n\cup_f e^{2n}$ będzie jej stożkiem
odwzorowania. Wybierzmy generatory
$a\in H^n(C_f;\Z)$ oraz $b\in H^{2n}(C_f;\Z)$.
Istnieje jednoznaczne $H(f)\in\Z$ takie, że
$a\smile a=H(f)b$. Liczbę $H(f)$ nazywamy
\emph{niezmiennikiem Hopfa} $f$ względem tych wyborów.
Przy ustalonej orientacji komórki górnej jej wartość
jest dobrze określona; zmiana generatora $b$ odwraca znak.
\end{definicja}

\begin{twierdzenie}[Kwadrat Whiteheada na \texorpdfstring{$S^2$}{S2}]
\label{thm:whitehead-square-v7}
Dla generatora $\iota_2\in\pi_2(S^2)$ i generatora
Hopfa $\eta\in\pi_3(S^2)\cong\Z$ zachodzi
$[\iota_2,\iota_2]=\pm2\eta$.
\end{twierdzenie}
\begin{proof}
Rozważmy CW-rozkład $S^2\times S^2$ z komórkami
w wymiarach $0,2,2,4$. Niech $x,y$ będą klasami
stopnia $2$ pochodzącymi z dwóch rzutowań, a $u$
generatorem stopnia $4$. W pierścieniu kohomologii
mamy $x^2=y^2=0$ i $x\smile y=u$,
po ustaleniu orientacji. Zatem
$(x+y)^2=2u$.

Niech $\omega:S^3\to S^2\vee S^2$ oznacza odwzorowanie
dołączającą komórkę czterowymiarową, a
$\nabla:S^2\vee S^2\to S^2$ odwzorowanie składającą oba
składniki przez identyczność. Stożek
$C_{\nabla\circ\omega}=S^2\cup_{\nabla\circ\omega}e^4$
jest skonstruowany z tej samej komórki górnej.
Z własności ilorazu dołączającego istnieje odwzorowanie
$F:S^2\times S^2\to C_{\nabla\circ\omega}$,
która na bukiecie jest równa $\nabla$, a na komórce
czterowymiarowej ma stopień $1$. Jeżeli $a,b$
są generatorami odpowiednio $H^2(C_{\nabla\circ\omega})$
i $H^4(C_{\nabla\circ\omega})$, to
$F^*a=x+y$ i $F^*b=\pm u$.
Zatem $F^*(a^2)=2u$, skąd, przez izomorfizm
$F^*:H^4(C_{\nabla\circ\omega})\to H^4(S^2\times S^2)$,
otrzymujemy $H(\nabla\circ\omega)=\pm2$.
Ponieważ niezmiennik Hopfa
$H:\pi_3(S^2)\to\Z$ jest izomorfizmem, wysyłającym
$\eta$ na $\pm1$, stąd $[\iota_2,\iota_2]=\pm2\eta$.
\end{proof}

\begin{przyklad}[Dwa przyklejenia tej samej komórki]
Rozważmy $S^2\vee S^4$ oraz
$\C P^2\simeq S^2\cup_{\eta}e^4$.
Obie przestrzenie mają $H_0=H_2=H_4=\Z$ i zerowe
pozostałe grupy homologii zredukowanej, ale
$H^*(S^2\vee S^4;\Z)$ ma zerowy kwadrat generatora
stopnia $2$, a w $H^*(\C P^2;\Z)$ jego kwadrat
jest generatorem $H^4$. Ta różnica jest ujęta
przez niezmiennik Hopfa odwzorowania dołączającego.
\end{przyklad}

\begin{przyklad}[Odwzorowanie z bukietu sfer nie przedłuża się na iloczyn]
\label{prz:whitehead-fold-product-v7}
Weźmy $p=q=2$, $X=S^2$ oraz
$\alpha=\beta=\iota_2$, czyli identyczność
na każdym z dwóch składników.
Ich suma bukietowa jest \emph{odwzorowaniem składającym}
\[
 \nabla:S^2\vee S^2\longrightarrow S^2,
 \qquad \nabla|_{S^2_1}=\nabla|_{S^2_2}=\id_{S^2}.
\]
Z Twierdzeń~\ref{thm:whitehead-square-v7}
i~\ref{thm:whitehead-product-v7} wynika, że
nie istnieje odwzorowanie
$g:S^2\times S^2\to S^2$
rozszerzające $\nabla$, ponieważ
$[\iota_2,\iota_2]=\pm2\eta\ne0$.

Sprawdźmy wynik \emph{niezależnym rachunkiem kohomologii}.
Niech $u\in H^2(S^2;\Z)$ będzie generatorem,
a $x=\operatorname{pr}_1^*u$,
$y=\operatorname{pr}_2^*u$
w $H^2(S^2\times S^2;\Z)$.
Jeżeli takie $g$ istniałoby,
to z ograniczeń do obu czynników
wynikałoby $g^*u=x+y$.
Ponieważ $u^2=0$, otrzymalibyśmy
\[
 0=g^*(u^2)=(g^*u)^2
 =(x+y)^2=x^2+2xy+y^2=2xy.
\]
Jednak $xy$ jest generatorem
$H^4(S^2\times S^2;\Z)\cong\Z$,
więc $2xy\ne0$.
Sprzeczność pokazuje także,
dlaczego współczynnik $2$ w kwadracie
Whiteheada ma znaczenie:
na poziomie samych grup
$H^2$ tej przeszkody nie widać.
\end{przyklad}

\begin{figure}[!htbp]
\centering
\begin{tikzpicture}[>=Stealth,font=\small]
 \node[draw=manifoldblue!70!black,rounded corners=4pt,
       fill=manifoldblue!8,minimum width=3.3cm,
       minimum height=1.15cm,align=center] (a) at (0,0)
       {$S^2\vee S^2$\\[-2pt] $\nabla=(\id,\id)$};
 \node[draw=manifoldgreen!65!black,rounded corners=4pt,
       fill=manifoldgreen!8,minimum width=3.3cm,
       minimum height=1.15cm,align=center] (b) at (6,0)
       {$S^2$};
 \node[draw=manifoldgold!85!black,rounded corners=4pt,
       fill=manifoldgold!10,minimum width=3.9cm,
       minimum height=1.25cm,align=center] (c) at (0,-2.2)
       {$S^2\times S^2$\\
        $(S^2\vee S^2)\cup_{\omega}D^4$};
 \draw[->,thick] (a)--(b) node[midway,above] {$\nabla$};
 \draw[hook->,thick] (a)--(c);
 \draw[->,densely dashed,manifoldred,thick]
       (c.east)--(b.south west)
       node[midway,below right] {$\nexists\,g$};
 \node[manifoldred,align=center] at (6,-2.1)
       {$[\iota_2,\iota_2]=\pm2\eta\ne0$};
\end{tikzpicture}
\caption{Górna komórka $D^4$ iloczynu $S^2\times S^2$
jest dołączona do bukietu $S^2\vee S^2$ przez
iloczyn Whiteheada. Złożenie $\nabla$ nie przedłuża
się na iloczyn: przeszkoda jest równa $\pm2\eta$.}
\label{fig:whitehead-product-obstruction-v7}
\end{figure}

\section{Podnoszenie przez wieżę Postnikowa i wyższe przeszkody}
\begin{twierdzenie}[Przeszkoda podniesienia]
\label{thm:postnikov-lift-v7}
Niech $B$ będzie jednospójnym CW-kompleksem, $A$
grupą abelową, $n\ge2$, a
$k:B\to K(A,n+1)$ odwzorowaniem. Oznaczmy przez
$p:E=\operatorname{hofib}(k)\to B$ odpowiadającą
fibrację. Dla CW-kompleksu $Z$ i ciągłej odwzorowania
$f:Z\to B$ istnieje odwzorowanie $\widetilde f:Z\to E$
ze złożeniem $p\circ\widetilde f=f$ wtedy i tylko
wtedy, gdy $f^*[k]=0$ w $H^{n+1}(Z;A)$.
\end{twierdzenie}
\begin{proof}
Z definicji $E$ punkt nad $b\in B$ jest parą
$(b,\gamma)$, gdzie $\gamma$ jest drogą od $k(b)$
do wyróżnionego punktu $K(A,n+1)$. Podniesienie
$\widetilde f(z)=(f(z),\gamma_z)$ jest ciągłe
dokładnie wtedy, gdy drogi $\gamma_z$ zależą
ciągle od $z$. To jest homotopia od $k\circ f$
do odwzorowania stałego. A zatem podniesienie istnieje
dokładnie wtedy, gdy $[k\circ f]=0$.
Na mocy twierdzenia o reprezentacji kohomologii
$[Z,K(A,n+1)]\cong H^{n+1}(Z;A)$, więc ten
warunek to $f^*[k]=0$.
\end{proof}

\begin{przyklad}[Pierwsza wyższa przeszkoda dla sfery]
Niech $u\in H^2(K(\Z,2);\Z)$ będzie generatorem,
a $Z$ dowolnym CW-kompleksem. Klasa
$c\in H^2(Z;\Z)$ reprezentuje odwzorowanie
$f:Z\to P_2S^2\simeq K(\Z,2)$. Istnienie jej
podniesienia do $P_3S^2$ wymaga i wystarcza,
aby $c\smile c=0$ w $H^4(Z;\Z)$, gdyż
$k_3=\pm u^2$. \emph{Nie} jest to samo co
istnienie odwzorowania $Z\to S^2$ indukującej $c$:
w dalszych etapach Postnikowa mogą pojawić się
nowe przeszkody. Jeżeli jednak $\dim Z\le3$,
przeszkoda w stopniu $4$ znika z powodów
wymiarowych, a dalsze również nie występują.
\end{przyklad}

\begin{przyklad}[Dlaczego generator $H^2(\C P^2;\Z)$ nie pochodzi z odwzorowania do $S^2$]
\label{prz:postnikow-cp2-nie-ma-liftu}
Oznaczmy przez $h\in H^2(\C P^2;\Z)$
dodatni generator. Z komórkowej struktury
$\C P^2=S^2\cup_\eta e^4$ oraz
rachunku iloczynu kubkowego wiemy, że
\[
 H^*(\C P^2;\Z)\cong\Z[h]/(h^3),
 \qquad h^2\ne0\in H^4(\C P^2;\Z).
\]
Twierdzenie~\ref{thm:rep-coho-v6}
przyporządkowuje $h$ odwzorowanie
$f:\C P^2\to K(\Z,2)=P_2S^2$.
Gdyby istniało podniesienie
$\widetilde f:\C P^2\to P_3S^2$,
Twierdzenie~\ref{thm:postnikov-lift-v7}
wymagałoby
$f^*k_3=\pm h^2=0$, co jest sprzecznością.
Zatem pierwsza nietrywialna przeszkoda
Postnikowa jest tutaj rzeczywiście niezerowa.

Można to zobaczyć także bez użycia wieży.
Gdyby istniało odwzorowanie $g:\C P^2\to S^2$
o $g^*u=h$, gdzie $u$ generuje
$H^2(S^2;\Z)$, to z naturalności
iloczynu kubkowego mielibyśmy
\[
 h^2=g^*(u\smile u)=g^*(0)=0,
\]
sprzecznie z powyższym rachunkiem.
W szczególności \emph{nie można} rozszerzyć
identyczności $S^2\to S^2$ przez
czterowymiarową komórkę dołączoną w $\C P^2$.
Przeszkodą jest klasa odwzorowania Hopfa
$\eta\in\pi_3(S^2)\cong\Z$,
która nie jest zerowa.
Dwa argumenty — z niezmiennika $k_3$
i z iloczynu kubkowego — wykrywają tę samą przeszkodę.
\end{przyklad}

\begin{uwaga}[Uwaga o wyborze podniesienia]
Jeżeli $f^*[k]=0$, to zbiór homotopii
zerujących $k\circ f$, \emph{z ustalonym $f$}
i z homotopiami włóknowymi, jest torsorem
nad $[Z,\Omega K(A,n+1)]\cong H^n(Z;A)$
(po wyborze jednego zerowania). Jeżeli dodatkowo
identyfikuje się podniesienia nad odwzorowaniemi $f$
zmieniającymi się przez homotopię, może pojawić się
działanie grupy pętli przestrzeni odwzorowań do $B$;
wtedy bez dodatkowych hipotez nie należy twierdzić,
że klasy wszystkich podniesień tworzą taki torsor.
\end{uwaga}

\section{Steenrod: od diagonalnej do działań na kohomologii}
\subsection*{Dlaczego zwykły iloczyn kubkowy nie wystarczy?}
Dla dowolnej przestrzeni $X$ z iloczynem kubkowym
znamy mnożenie $H^p(X;\F_2)\otimes H^q(X;\F_2)
\to H^{p+q}(X;\F_2)$, gdzie $\F_2=\Z/2\Z$.
Jednak nie każda operacja kohomologiczna jest
uzyskiwana przez pomnożenie dwóch niezależnie
wybranych klas. Interesują nas \emph{naturalne}
odwzorowania zwiększające stopień
$H^n(X;\F_2)\to H^{n+i}(X;\F_2)$,
zdefiniowane dla wszystkich przestrzeni $X$.

\begin{definicja}[Wyższe diagonalne i iloczyny $\smile_r$]
Niech $C_*(X;\F_2)$ będzie unormowanym singularnym
kompleksem łańcuchowym, $\delta$ jego koróżniczką
na $C^*(X;\F_2)$, a $T(a\otimes b)=b\otimes a$
zamianą czynników tensora (nad $\F_2$ znaki
Koszula znikają). \emph{Wyższymi przybliżeniami diagonalnej}
nazywamy naturalne odwzorowania stopnia $r$
\(\Delta_r:C_m(X;\F_2)\to
(C_*(X;\F_2)\otimes C_*(X;\F_2))_{m+r}\),
$r\ge0$, dla których $\Delta_0$ jest diagonalną
Alexandra--Whitneya i
\[
\partial\Delta_r+\Delta_r\partial=(1+T)\Delta_{r-1}
\quad(r\ge1).
\]
Dla jednorodnych kołańcuchów $a,b$ definiujemy
$\smile_r$ przez
$(a\smile_r b)(c)=(a\otimes b)(\Delta_r c)$,
a $\smile_{-1}$ określamy jako zero.
Wtedy $a\smile_r b$ ma stopień
$\deg a+\deg b-r$.
\end{definicja}

\begin{lemat}[Istnienie i tożsamość koróżniczki]
\label{lem:cup-i-v7}
Istnieją naturalne operatory $\Delta_r$. Korzystamy
z klasycznej unormowanej konstrukcji Steenroda; można je
wybrać tak, by $\smile_0$ zgadzał się z iloczynem
Alexandra--Whitneya, a $a\smile_{\deg a}a=a$
dla kołańcucha o wartościach w $\F_2$ w odpowiednim
stopniu. Dla jednorodnych kołańcuchów zachodzi
\[
\delta(a\smile_r b)=(\delta a)\smile_r b
+a\smile_r(\delta b)+a\smile_{r-1}b+b\smile_{r-1}a.
\]
\end{lemat}
\begin{proof}
Opiszmy konstrukcję komórkową, nie tylko odwołanie
do istnienia. Na singularnym sympleksie $\sigma:\Delta^m\to X$
operatory wyznacza się na standardowym sympleksie
$\Delta^m$, a następnie przenosi przez $\sigma$.
Dla $r=0$ bierzemy zwykły wzór Alexandra--Whitneya:
\(\Delta_0[v_0\dots v_m]=
\sum_{i=0}^m[v_0\dots v_i]\otimes[v_i\dots v_m]\).

Załóżmy, że $\Delta_0,\ldots,\Delta_{r-1}$
zostały już określone naturalnie. Dla każdego
$\Delta^m$ chcemy rozwiązać równanie
$\partial\Delta_r+\Delta_r\partial=(1+T)\Delta_{r-1}$.
Na brzegu sympleksu $\Delta_r$ jest już znana
z indukcji po wymiarze $m$. Żądany brzeg
$\Delta_r[\Delta^m]$ to suma
$(1+T)\Delta_{r-1}[\Delta^m]
+\Delta_r(\partial[\Delta^m])$.
Jest ona cyklem: stosując $\partial$ i hipotezę
indukcyjną otrzymujemy $(1+T)^2\Delta_{r-2}=0$
ponieważ $T^2=1$ w charakterystyce $2$.
Grupa docelowa jest tensorowym kompleksem
łańcuchowym standardowego sympleksu z sobą;
jest on acykliczny w dodatnich stopniach, gdyż
$\Delta^m$ jest ściągalny, a łańcuchy są wolne.
Żądany cykl leży w stopniu $m+r-1>0$
poza krańcowym przypadkiem $m=0,r=1$,
który ustala się bezpośrednio. Można go więc
wypełnić, wybierając $\Delta_r[\Delta^m]$.
To daje jednocześnie indukcję po $m$ i $r$.
Przenoszenie operatorów ze standardowych sympleksów przez
odwzorowania singularne zapewnia naturalność względem
ciągłych odwzorowań. \emph{Trzeba przy tym zaznaczyć
istotne ograniczenie tego argumentu:} sama acykliczność
modeli daje rozwiązania relacji homotopii, lecz nie dowodzi
jeszcze warunku normalizacji
$a\smile_n a=a$ dla dowolnego kołańcucha stopnia $n$.
Do uzyskania tej normalizacji wybieramy standardowe
\emph{unormowane wyższe diagonalne Steenroda}; ich
jawna konstrukcja wymaga dodatkowych warunków na wypełnienia
(np. wzorów przedziałowych na uporządkowanych sympleksach).
W dalszym ciągu normalizacja jest zatem przywoływaną
własnością tej konkretnej klasycznej konstrukcji, a nie
wnioskiem z samego rozwiązywania równania na $\Delta_r$.

Ostatnią tożsamość otrzymujemy, wartościując relację
na $\Delta_r$ przez $a\otimes b$. Człony
z $\partial$ dają koróżniczki obu czynników,
a człon $(1+T)\Delta_{r-1}$ daje
$a\smile_{r-1}b+b\smile_{r-1}a$.
\end{proof}

\begin{definicja}[Kwadraty Steenroda]
Niech $X$ będzie przestrzenią topologiczną,
$n\ge0$, $i\ge0$. Dla kocyklu
$a\in Z^n(X;\F_2)$ definiujemy
\[
\operatorname{Sq}^i[a]=
\begin{cases}
[a\smile_{n-i}a]\in H^{n+i}(X;\F_2),&0\le i\le n,\\
0,&i>n.
\end{cases}
\]
Operacje $\operatorname{Sq}^i$ nazywamy
\emph{kwadratami Steenroda}.
\end{definicja}

\begin{twierdzenie}[Poprawność definicji i własności podstawowe]
\label{thm:steenrod-basic-v7}
Dla każdej przestrzeni $X$ operacje
$\operatorname{Sq}^i:H^n(X;\F_2)\to H^{n+i}(X;\F_2)$
są dobrze określonymi naturalnymi homomorfizmami
przestrzeni wektorowych nad $\F_2$. Ponadto:
$\operatorname{Sq}^0=\id$,
$\operatorname{Sq}^n(x)=x\smile x$ dla $x\in H^n$,
a $\operatorname{Sq}^i(x)=0$ dla $i>\deg x$.
\end{twierdzenie}
\begin{proof}
Oznaczmy $r=n-i$. Dla kocyklu $a$ tożsamość
Lematu~\ref{lem:cup-i-v7} daje
$\delta(a\smile_r a)=2a\smile_{r-1}a=0$.
Zatem przedstawiona formuła wyznacza klasę
kohomologii. Niech $a'=a+\delta b$, gdzie
$b\in C^{n-1}(X;\F_2)$. Po rozwinięciu
$a'\smile_r a'$ różnica w stosunku do
$a\smile_r a$ jest sumą
\[
a\smile_r\delta b+\delta b\smile_r a
+\delta b\smile_r\delta b.
\]
Jest to koróżniczka kołańcucha
\(a\smile_r b+b\smile_r a
+b\smile_r\delta b+b\smile_{r-1}b\):
pierwsze dwa składniki otrzymujemy z tożsamości
dla $a\smile_r b+b\smile_r a$,
a pozostałe z tożsamości dla
$b\smile_r\delta b+b\smile_{r-1}b$.
Gdy $r=0$, składnik z $\smile_{-1}$ jest zerowy.
Zatem wynik nie zależy od reprezentanta.

Dla kocykli $a,c\in Z^n$ wyrazy mieszane
$a\smile_r c+c\smile_r a$ są koróżniczką
$a\smile_{r+1}c$. Wobec tego
$\operatorname{Sq}^i([a]+[c])=
\operatorname{Sq}^i[a]+\operatorname{Sq}^i[c]$.
Naturalność wynika z naturalności $\Delta_r$.
Normalizacja $a\smile_n a=a$ daje
$\operatorname{Sq}^0=\id$, zaś dla $i=n$
mamy $r=0$, czyli zwykły iloczyn kubkowy.
Przypadek $i>n$ jest częścią definicji.
\end{proof}

\begin{twierdzenie}[Wzór Cartana i relacje Adema]
\label{thm:steenrod-cartan-v7}
Dla $x\in H^p(X;\F_2)$, $y\in H^q(X;\F_2)$
oraz $k\ge0$ zachodzi \emph{wzór Cartana}
\[
\operatorname{Sq}^k(x\smile y)=
\sum_{i+j=k}\operatorname{Sq}^i(x)\smile\operatorname{Sq}^j(y).
\]
Ponadto dla $0<a<2b$ zachodzą \emph{relacje Adema}
\[
\operatorname{Sq}^a\operatorname{Sq}^b
=\sum_{j=0}^{\lfloor a/2\rfloor}
\binom{b-j-1}{a-2j}\operatorname{Sq}^{a+b-j}\operatorname{Sq}^j,
\]
gdzie współczynniki dwumianowe liczymy modulo $2$,
a iloczyn operatorów oznacza ich złożenie.
\end{twierdzenie}
\begin{uwaga}[Szkic i status dowodów]
Wzór Cartana wynika z porównania wyższych diagonalnych
dla tensora kompleksów z tensorami diagonalnych
obu czynników: składniki stopnia $k$ powstają ze
wszystkich rozkładów $k=i+j$. Zgodność wymaga jednak
jawnych homotopii między diagonalnymi ekwiwariantnymi.
Relacje Adema wymagają jeszcze wyższych homotopii
dla \emph{iterowanych} diagonalnych. W obecnej iteracji
podajemy oba standardowe twierdzenia i ich zastosowania,
ale \emph{nie deklarujemy pełnego dowodu Cartana ani
Adema}. Stanowią one odrębny etap konstrukcyjny
(teoria spójnych diagonalnych i algebry Steenroda).
Natomiast dowody istnienia kwadratów, ich naturalności,
addytywności, niestabilności oraz górnego kwadratu
zostały podane wyżej.
\end{uwaga}

\section{Przykłady operacji Steenroda}

\begin{przyklad}[Pierwszy kwadrat na trójkącie]
Niech $K$ będzie kompleksem symplicjalnym i
$a\in Z^1(K;\F_2)$. Dla uporządkowanej
zorientowanej $2$-komórki $[v_0,v_1,v_2]$
zwykły iloczyn kubkowy jest dany przez
$\bigl(a\smile a\bigr)[v_0,v_1,v_2]
=a[v_0,v_1]a[v_1,v_2]$. Zatem
$\operatorname{Sq}^1[a]=[a\smile a]$.
To wzór wykonawczy: wystarczy obliczyć wartości
kocyklu na krawędziach. Dla $a$ stopnia $1$
operacje wyższe niż $\operatorname{Sq}^1$ znikają
na mocy niestabilności.
\end{przyklad}

\begin{przyklad}[\texorpdfstring{$\R P^\infty$}{RP infinity}]
Niech $t\in H^1(\R P^\infty;\F_2)$ będzie generatorem.
Mamy $H^*(\R P^\infty;\F_2)\cong\F_2[t]$.
Z własności podstawowych wynika
$\operatorname{Sq}^0(t)=t$,
$\operatorname{Sq}^1(t)=t^2$ i
$\operatorname{Sq}^i(t)=0$ dla $i>1$.
Stosując wzór Cartana do $t^m$, otrzymujemy
\[
\operatorname{Sq}^i(t^m)=\binom mi t^{m+i}.
\]
W $\R P^N$ ten sam wzór obowiązuje, jeśli
$m+i\le N$; w większych stopniach prawa strona
znika ze względów wymiarowych.
\end{przyklad}

\begin{przyklad}[\texorpdfstring{$\C P^\infty$}{CP infinity}]
Niech $u\in H^2(\C P^\infty;\F_2)$ będzie redukcją
modulo $2$ uniwersalnej klasy Cherna. Ponieważ
grupy w nieparzystych stopniach są zerowe,
$\operatorname{Sq}^1(u)=0$, a z własności górnego
kwadratu $\operatorname{Sq}^2(u)=u^2$.
Wzór Cartana daje
$\operatorname{Sq}^{2j}(u^m)=\binom mj u^{m+j}$,
a $\operatorname{Sq}^{2j+1}(u^m)=0$.
\end{przyklad}

\begin{przyklad}[Kwadraty Steenroda różnią się od potęgowania]
\label{prz:steenrod-rp5-obliczenia-v7}
Niech $t\in H^1(\R P^5;\F_2)$ będzie standardowym
generatorem. Pierścień kohomologii jest
$H^*(\R P^5;\F_2)=\F_2[t]/(t^6)$.
Dla klasy $x=t^3\in H^3(\R P^5;\F_2)$
formuła dwumianowa z poprzedniego przykładu daje
\[
 \operatorname{Sq}^1(x)=\binom31t^4=t^4\ne0,
 \qquad
 \operatorname{Sq}^2(x)=\binom32t^5=t^5\ne0.
\]
Tymczasem \emph{kwadrat kubkowy} $x^2=t^6$
jest zerem ze względów wymiarowych.
Operacje $\operatorname{Sq}^1$ i
$\operatorname{Sq}^2$ wykrywają zatem
niezerowe klasy, choć zwykłe potęgowanie
$x\smile x$ nie wykrywa niczego.

Jeszcze jedno sprawdzenie dotyczy relacji Adema.
Dla każdego $m\ge0$ w $\R P^\infty$
obliczamy
\[
 \operatorname{Sq}^1(\operatorname{Sq}^1(t^m))
 =\binom m1\binom{m+1}1t^{m+2}
 =m(m+1)t^{m+2}=0
\]
w charakterystyce $2$. Zgadza się to
z relacją $\operatorname{Sq}^1\operatorname{Sq}^1=0$.
Ten rachunek \emph{weryfikuje przykład},
ale nie zastępuje ogólnego dowodu relacji Adema
dla dowolnej przestrzeni.
\end{przyklad}

\begin{przyklad}[Zastosowanie do cofnięć]
Jeśli $f:X\to\R P^\infty$ klasyfikuje rzeczywistą
wiązkę liniową $L$ i $w_1(L)=f^*t$, to naturalność
kwadratów Steenroda daje
$\operatorname{Sq}^1(w_1(L))=w_1(L)^2$.
Dla przestrzeni o większej liczbie generatorów
operacje potrafią odróżnić odwzorowania, które na poziomie
addytywnym dają te same grupy kohomologii.
\end{przyklad}

\section{Przygotowanie do dalszej teorii}
\begin{uwaga}[Przypadek niejednospójny]
Jeżeli $\pi_1(X)$ nie jest zerowa, jej działanie
na $\pi_n(X)$ dla $n\ge2$ musi wejść do danych
wieży Postnikowa. Odpowiednie niezmienniki $k$
należą do grup kohomologii z \emph{lokalnymi
współczynnikami}. W szczególności nie wolno
zastępować ich bez uzasadnienia grupami
$H^{n+1}(P_{n-1}X;\pi_nX)$ ze stałymi współczynnikami.
\end{uwaga}

\begin{uwaga}[Nieskończone wieże i zbieżność]
Skończone etapy Postnikowa opisują kolejno coraz
więcej grup homotopii. Przejście do granicy
homotopijnej całej wieży jest odrębnym zagadnieniem:
w ogólnej sytuacji występują warunki kompletności
i możliwe człony $\varprojlim^1$. W tym rozdziale
nie utożsamiamy automatycznie $X$ ze zwykłą
granicą odwrotną $\varprojlim P_nX$.
\end{uwaga}

\begin{uwaga}[Miejsce geometrii algebraicznej]
W późniejszym połączeniu z geometrią algebraiczną
ciągi spektralne Serre'a będą punktem odniesienia
dla ciągu Leraya, a operacje Steenroda dla analogicznych
operacji w kohomologii motywicznej. Same
konstrukcje analityczne, etalne i snopowe wymagają
jednak osobnych hipotez i nie są utożsamiane z
operacjami zdefiniowanymi dla topologicznych
CW-kompleksów.
\end{uwaga}


% Iteracja VIII, rozdział główny 15. Do umieszczenia przed rozdziałem o stabilizacji.
